\subsection{Определение жадных алгоритмов}

Это алгоритмы, заключающиеся в принятии локально оптимальных решений на каждом этапе, допуская, что конечное решение
также окажется оптимальным

\subsection{Процедура построения}

В ходе построения дерева мы разбиваем вашу выборку на две части.
Для этого представим, что у нас есть $L$ и $R$ — левое и правое поддерево объектов, тогда вспомним, что
разбиение в узле строится по критерию вида $\mathbb{I} [x_i^j < b]$, $b$ -- порог, $x_i^j$ -- $j$-й признак у $i$-го
объекта.
Осталось понять, как выбрать $j$ и $b$.
Для этого будем перебирать все возможные признаки $j$, а для каждого признака перебирать все его возможные значения
$x^j$ в выборке, тогда мы можем построить разбиение для каждого из вариантов, осталось выбрать оптимальный.
Для этого используют функционал следующего вида: $f(X) = \frac{\abs{L}}{\abs{X}} H(L) + \frac{\abs{R}}{\abs{X}} H(R)$

\subsection{Информационные критерии, энтропия, индекс Джини}

Функция $H$ называется критерием информативности (information criteria).
Есть разные способы задать $H$  для задачи классификации.
Посмотрим сначала на примере бинарной классификации.
Далее $p_i$ -- нормированная частота встречаемости $i$-ого класса в множестве $R$

Наиболее простым является Misclassification criteria: количество ошибок классификации (предсказываем доминирующий
класс и смотрим сколько объектов классифицировали неправильно)
\[ H(R) = 1 - \max (p_0, p_1) \]
Не очень хороший способ: например для многоклассовой классификации дерево может хорошо отделять несколько каких-то
классов от остальных, но этот критерий нам этого не покажет

Энтропия: чем больше энтропия, тем больше хаоса
\[ H(R) = - p_0 \log p_0 - p_1 \log p_1 = - p_0 \log p_0 - (1 - p_0) \log (1 - p_0) \]

Индекс Джини:
\[ H(X_m) = \sum_{k=1}^K p_k (1 - p_k)\]

Критерий Джини допускает и следующую интерпретацию: $H(X_m)$ равно математическому ожиданию числа неправильно
классифицированных объектов в случае, если мы будем приписывать им случайные метки из дискретного распределения,
заданного вероятностями $(p_1, \dots, p_k)$.

\subsection{Ответ в листе}

Если в лист дерева попали объекты разных классов и нет задачи получить вектор вероятностей, то просто наиболее
часто встречаемый в листе класс.
Если надо еще получить вектор вероятностей для каждого класса, то вернем вектор из частот встречаемости в листе.
Если подробнее, то надо расписать лагранжиан

\subsection{Обработка пропусков}

Одна из приятных особенностей деревьев -- это способность обрабатывать пропуски в данных.
Разберёмся, что при этом происходит на этапе обучения и на этапе применения дерева.

Пусть у нас есть некоторый признак $x_i$, значение которого пропущено у некоторых объектов.
Как обычно, обозначим через $X_m$ множество объектов, пришедших в рассматриваемую вершину, а через $V_m$
-- подмножество $X_m$, состоящее из объектов с пропущенным значением $x_i$.
В момент выбора сплитов по этому признаку мы будем просто игнорировать объекты из $V_m$, а когда сплит выбран, мы
отправим их в оба поддерева.
При этом логично присвоить им веса: $\frac{\abs{X_l}}{\abs{X_m}}$ для левого поддерева и
$\frac{\abs{X_r}}{\abs{X_m}}$ для правого.
Веса будут учитываться как коэффициенты при $L(y_i,c)$ в формуле информативности.

Теперь рассмотрим этап применения дерева.
Допустим, в вершину, где сплит идёт по $i$-му признаку, пришёл объект $x_0$ с пропущенным значением этого признака.
Предлагается отправить его в каждую из дальнейших веток и получить по ним предсказания $y^l$ и $y^r$.
Эти предсказания мы усредним с весами $\frac{\abs{X_l}}{\abs{X_m}}$ и $\frac{\abs{X_r}}{\abs{X_m}}$
(которые мы запомнили в ходе обучения):
\[ y = \frac{\abs{X_l}}{\abs{X_m}} y^l + \frac{\abs{X_r}}{\abs{X_m}} y^r \]

Для задачи регрессии это сразу даст нам таргет, а в задаче бинарной классификации -- оценку вероятности класса 1.
Заметим, что если речь идёт о категориальном признаке, может оказаться хорошей идеей ввести дополнительное значение
«пропущено» для категориального признака и дальше работать с пропусками, как с обычным значением.
Особенно это актуально в ситуациях, когда пропуски имеют системный характер и их наличие несёт в себе определённую
информацию

\subsection{Бинаризация}

Идея состоит в выборе нескольких порогов для одного признака, вместо одного
Это настраивается, например, by Border Count Hyperparameter in Catboost (defaults to 254 = 256 - inf)
То есть если подпространство исходного пространства признаков, то в соответствии с листом дерева предсказание
будет иметь вид $y = \sum_j w_j [x \in J_j] $

\subsection{Обрезка (pruning), варианты C4.5 и CART}

Мы уже упоминали выше, что деревья легко переобучаются и процесс ветвления надо в какой-то момент останавливать.
Для этого есть разные критерии, обычно используются все сразу

\begin{itemize}
    \item Ограничение по максимальной глубине дерева;
    \item ограничение на минимальное количество объектов в листе;
    \item ограничение на максимальное количество листьев в дереве;
    \item требование, чтобы функционал качества BranchBranch при делении текущей подвыборки на две улучшался не
    менее чем на $s$ процентов.
\end{itemize}

Делать это можно на разных этапах работы алгоритма, что не меняет сути, но имеет разные устоявшиеся названия

\begin{itemize}
    \item Можно проверять критерии прямо во время построения дерева, такой способ называется pre-pruning или early
    stopping;
    \item а можно построить дерево жадно без ограничений, а затем провести стрижку (pruning), то есть удалить
    некоторые вершины из дерева так, чтобы итоговое качество упало не сильно, но дерево начало подходить под условия
    регуляризации.
    При этом качество стоит измерять на отдельной, отложенной выборке.
\end{itemize}

Стандарт C4.5
\begin{itemize}
    \item Критерий нормализованной энтропии
    \item Остановка зависит от размера листа.
    \item Включает обрезку
\end{itemize}

Стандарт CART
\begin{itemize}
    \item Критерии Джини
    \item Сокращение затрат и сложности
    \item Замещающие предикаты для отсутствующих данных
\end{itemize}

\subsection{Гиперпараметры (минимум 4 основных)}

\begin{itemize}
    \item Максимальная глубина
    \item Минимальное число элементов в листьях (и выбор условия в листе)
    \item Минимальное разделение образцов
    \item Минимальное количество примесей
\end{itemize}